\begin{comment}

The protonation of a flexible molecule cannot generally be separated from its conformation. Protonation changes charge distribution, hydrogen bonding, structure, reactivity, and optical response, while conformational changes alter the solvation and proton affinity of individual sites.\cite{manallackPKaDistributionDrugs2007,martinLetsNotForget2009,baxGettingChemistryRight2017,maBiomedicalApplicationsSupramolecular2015,muruganRoleProtonationState2014,handeRecentAdvancesSmall2022} A molecule containing several titratable and tautomerizable sites must therefore be described by two coupled ensembles: an ensemble of conformations and, within each conformation, an ensemble of protonation tautomers. Their joint populations depend on pH, solvent, and the relative free energies of both protonation and structural rearrangement.\cite{gunnerStandardStateFree2020,isikOverviewSAMPL6PKa2021,nagyTheoreticalCalculationsConformational2014}

A measured \pka{} reports a change in the total proton occupancy of this combined ensemble. It does not, by itself, reveal the molar fractions of the underlying tautomers or the conformations from which proton exchange occurs.\cite{ullmannRelationsProtonationConstants2003,bombardaPHDependentPKaValues2010} This distinction becomes particularly important when conformation and acidity are strongly coupled. A conformer with a small equilibrium population can contribute disproportionately to the apparent acidity if its proton affinity differs sufficiently from that of the dominant conformer. Conversely, protonation can redistribute the conformational ensemble. Consequently, even structures conventionally assigned the same conformational label may not be chemically equivalent for acid--base equilibria.

Biliverdin (BV) is an archetypal example of this coupled-ensemble problem. Its tetrapyrrole core contains four pyrrole rings, A to D, connected by flexible methine bridges. The four nitrogens can carry protons, while the terminal lactam groups can additionally form keto--enol tautomers (Figure~\ref{fig:micro}). In the cation, all four nitrogens are protonated; removal of one proton produces four principal neutral tautomers, distinguished by the nitrogen from which the proton is lost. Two propionic acid side chains introduce two additional titratable groups.

BV is often described by a single configurational label, such as the helical \(ZZZsss\) form. In solution, however, \(ZZZsss\) denotes a broad basin containing many interconverting structures with different helicities, hydrogen-bonding patterns, solvent exposure, and relative orientations of the terminal rings. Opening one of the methine bridges accesses additional basins that can have substantially different acidities. Thus, assigning one structure and one protonation state to BV suppresses precisely the fluctuations that may control its observed acid--base behavior.

This coupling is also relevant to BV photochemistry. BV is a central intermediate in heme catabolism and the chromophore of bacteriophytochromes and near-infrared fluorescent proteins.\cite{tenhunen_microsomal_1969,chernovNearInfraredFluorescentProteins2017,tangRedEdgeBilinBinding2021} Its conformation, protonation, and hydrogen-bonding interactions influence spectral tuning and the initiation of photochemical signaling.\cite{modiProtonationBiliverdinIXa2019a,shankarUltrafastRemotecontrolledProtonation2025,songRedGreenCyanobacteriochrome2015,tohProtontransferHydrogenbondInteractions2010a,yangUltrafastProtoncoupledIsomerization2022a} The central question in solution is therefore not simply which nitrogen is protonated. It is how BV distributes its population across conformations and neutral tautomers, and how structural opening changes both the overall \pka{} and the tautomeric composition.

Experiment already indicates that BV acidity is highly conformation-dependent. Spectrophotometric titrations of free BV and related biliverdins give apparent \pka{} values spanning several units (Table~\ref{tab:literature}). In aqueous buffer, the Soret band of BV shifts from 379 to 354~nm with an apparent \pka{} of 5.7. The BV diamide gives approximately the same value, placing the transition on the tetrapyrrole core rather than the propionic acids.\cite{shangBiliverdinAmidesReveal2010} In alcohol--water mixtures, BV esters are less basic, with reported core \pka{} values between 3.6 and 4.3.\cite{holzwarthPhytochromeModelsII1978,kroisGeometryBasicityBilatrienes1991,bariInterplayBasicityConformation1992} More strikingly, synthetic biliverdins constrained into progressively extended conformations have \pka{} values of 6.1, 6.4, and 7.9, whereas a biliverdin locked into a compact helix has a \pka{} of only 0.4.\cite{bariInterplayBasicityConformation1992,kroisGeometryBasicityBilatrienes1991}

% Keep Table~\ref{tab:literature} here.

These measurements demonstrate a strong relationship between conformation and acidity, but they do not isolate its molecular origin. The experiments differ in solvent, side-chain derivatization, ionic composition, aggregation state, and conformational restriction. The aqueous measurement was performed on the free acid, whose propionic side chains titrate in the same pH range,\cite{lightnerAcidDissociationConstants1996} whereas most measurements in mixed solvents used esterified derivatives. Aggregation\cite{haidlNewAspectsProtonation1985} and protonation-induced structural changes\cite{kroisGeometryBasicityBilatrienes1991} can further shift the apparent \pka{}. Moreover, an absorption titration reports the total ensemble response and does not determine the molar fractions of the four neutral tautomers. Although one tautomer is commonly drawn as the dominant form, \textsuperscript{13}C NMR indicates rapid exchange between at least the two central N--H tautomers.\cite{petrier13CNuclearMagnetic1983,falkBeitrgeZurChemie1985}

Resolving this problem computationally requires the protonation and conformational coordinates to be treated on the same thermodynamic footing. Absolute proton-binding free energies are notoriously difficult to calculate because they combine electronic-structure errors, solvation free energies, standard-state conventions, and the long-range reorganization of water. Even for small aqueous acids, density-functional schemes commonly retain uncertainties of one to two \pka{} units.\cite{chengRedoxPotentialsAcidity2014} For BV, this difficulty is compounded by the number of possible tautomers and conformations. The acidity of each structure depends on where the remaining protons reside, while a conformer that is rare in the protonated ensemble may become important after deprotonation. Such conformers cannot be discarded solely because their initial equilibrium populations are small. Unbiased simulations of carbonic acid, for example, show that a conformer representing only about 12\% of the ensemble accounts for most proton-release events.\cite{zhaoInitioDeepNeural2025}

An explicit-water calculation must additionally distinguish local proton motion from completed deprotonation. Deprotonation is complete only when the excess-proton defect leaves the molecule and reaches bulk water. The defect migrates by Grotthuss structural diffusion through transient hydrogen-bonded water wires, so no single identifiable hydrogen atom travels from the solute to the bulk.\cite{agmonGrotthussMechanism1995,agmonProtonsHydroxideIons2016,cukiermanTuGrotthussOther2006,hassanaliProtonTransferWater2013,miyakeGrotthussMechanismsProton2016,popovSearchGrotthussMechanism2023} A trajectory may instead return the proton to its original site, transfer it to another solute site, trap it near the molecular surface, or recombine it with a counterion. These outcomes can appear similar along a local coordination coordinate but represent different chemical processes.

Tracking and biasing this delocalized defect normally requires a dedicated reactive model and system-specific parameterization.\cite{liUnderstandingTrackingExcess2020,grifoniMicroscopicDescriptionAcidbase2019} Extending such models to BV is especially demanding because the solvated system approaches \(10^3\) quantum-mechanical atoms and must be propagated long enough for proton transfer, solvent reorganization, and conformational motion. Conventional density-functional molecular dynamics is not practical at this scale.

We therefore use divide-and-conquer density-functional tight binding (DC-DFTB), which combines a quantum-mechanical description of bond breaking and proton transfer with near-linear scaling through spatial decomposition.\cite{gausDensityFunctionalTight2014,spiegelmanDensityfunctionalTightbindingBasic2020,nakaiDivideandConquerTypeDensityFunctionalTightBinding2016,nishimuraHierarchicalParallelizationDivideandconquer2020} DFTB proton-transfer calculations have been benchmarked on small aqueous acids and amines, where they reproduce the relevant acid--base energetics with sufficient accuracy for enhanced-sampling simulations.\cite{Tummanapelli2014,Tummanapelli2015GSH,Tummanapelli2015Amino,Sakti2018} Its scaling makes it possible to treat BV and its surrounding water within the same reactive quantum-mechanical calculation rather than dividing them into predefined reactive and nonreactive regions.

In this work, we combine reactive DC-DFTB metadynamics in explicit water with an analysis that follows the excess-proton charge defect through the solvent.\cite{barducciMetadynamics2011,nishizawaThreePillarsAchieving2016} A coordination-number collective variable promotes proton release without specifying the acceptor or final proton location,\cite{Tummanapelli2015Polyprotic,Tummanapelli2015Amino,Sakti2018} while postprocessing identifies whether the defect reaches bulk water and separates completed deprotonation from return, intramolecular transfer, trapping, or counterion recombination. We calculate one absolute acid--base equilibrium directly. Umbrella-sampling free energies for proton transfer between neighboring nitrogens then determine the relative populations of the remaining neutral tautomers, while torsional umbrella sampling quantifies how opening the helix changes their collective acidity. Together, these calculations form thermodynamic cycles that connect protonation, tautomerization, and conformation without requiring a separate absolute deprotonation calculation for every state.

We test the components of the approach on three smaller systems: histidine, which contains two competing sites with known microscopic acidities;\cite{kyte_structure_2006,tanokura1HNMRStudyTautomerism1983} propionic acid, whose \emph{syn} and \emph{anti} conformers provide a closed cycle between deprotonation and rotation;\cite{macoas_internal_2005} and malonaldehyde, a symmetric reference for intramolecular proton transfer.\cite{baughcumMicrowaveSpectroscopicStudy1981,wangFulldimensionalQuantumCalculations2008}

For BV, the compact helical ensemble has a macroscopic core \pka{} of 4.0, and all four neutral tautomers coexist without a single form dominating. Opening the helix at ring D makes the core substantially more basic, raising its macroscopic \pka{} to 6.4, and redistributes the neutral ensemble so that the ring-B and ring-C tautomers dominate. These results show that BV acidity is a property of a coupled conformational and protonation ensemble---not of a single optimized structure or a uniquely assigned nitrogen---and provide a molecular explanation for the strong conformational dependence observed experimentally.

\end{comment}
